%%%%%%%%%%%%%%%%%%%- HXN
\def\sode{12}
\def\tendethi{ĐỀ PHÁT TRIỂN MINH HOẠ 2025}
\begin{dethi}
 {\tendethi}
\end{dethi}
\caulc
\Opensolutionfile{ans}[ans/ans-HXN-\sode-T]
\begin{ex}%Câu 1
 Tìm tập xác định của hàm số $ y=\log_2\left(x-3\right)$.
 \choice
 {$\mathscr{D}=\left(-\infty;3\right)$}
 {$\mathscr{D}=\mathbb{R}$}
 {\True $\mathscr{D}=\left(3;+\infty\right)$}
 {$\mathscr{D}=\left[3;+\infty\right)$}
\end{ex}
\begin{ex}%Câu 2
 Số đường tiệm cận của đồ thị hàm số $y=\dfrac{3x-4}{x-1}$ bằng
 \choice
 {\True $2$}
 {$3$}
 {$1$}
 {$0$}
\end{ex}
\begin{ex}%Câu 3
 Độ $pH$ của một dung dịch được tính theo công thức $ pH=-\log\left[H^+\right]$ với $\left[H^+\right]$ là nồng độ ion hydrogen của dung dịch đó. Độ $ pH$ của một loại sữa có $\left[H^+\right]=10^{-6{,}8}$ là bao nhiêu?
 \choice
 {$-6{,}8$}
 {$68$}
 {\True $6{,}8$}
 {$0{,}68$}
\end{ex}
\begin{ex}%Câu 4
 Cho hình lăng trụ đứng $ABC.A'B'C'$ có $BB'=a$ , đáy $ABC$ là tam giác vuông cân tại $ B$ và $BA=BC=a$ . Tính thể tích $ V$ của khối lăng trụ đã cho.
 \choice
 {$V=a^3$}
 {$V=\dfrac{a^3}{3}$}
 {$V=\dfrac{a^3}{6}$}
 {\True $V=\dfrac{a^3}{2}$}
\end{ex}
\begin{ex}%Câu 5
 Trong không gian với hệ tọa độ $ Oxyz$, cho mặt phẳng $\left(\alpha\right)\colon x-2y+2z-3=0.$ Điểm nào sau đây thuộc mặt phẳng $\left(\alpha\right)$?
 \choice
 {$M\left(2;0;1\right)$}
 {$Q\left(2;1;1\right)$}
 {$P\left(2;-1;1\right)$}
 {\True $N\left(1;0;1\right)$}
\end{ex}
\begin{ex}%Câu 6
 Cho hình hộp $ ABCD.A'B'C'D'$. Vectơ $\overrightarrow{v}=\overrightarrow{B'A'}+\overrightarrow{B'C'}+\overrightarrow{B'B}$ bằng vectơ nào dưới đây?
 \choice
 {$\overrightarrow{DB'}$}
{$\overrightarrow{B'D'}$}
{$\overrightarrow{BD'}$}
{\True $\overrightarrow{B'D}$}
\end{ex}
\begin{ex}%Câu 7
Họ nguyên hàm của hàm số $ f(x)=3x^2+\sin x$ là 
\choice
{$x^3+\cos x+C$}
{$x^3+\sin x+C$}
{\True $x^3-\cos x+C$}
{$3x^3-\sin x+C$}
\end{ex}
\begin{ex}%Câu 8
Trong không gian $ Oxyz$, phương trình nào sau đây là phương trình của mặt cầu có tâm $ I\left(7;6;-5\right)$ và bán kính $ 9$?
\choice
{$\left(x+7\right)^2+\left(y+6\right)^2+\left(z-5\right)^2=81$}
{$\left(x+7\right)^2+\left(y+6\right)^2+\left(z-5\right)^2=9$}
{\True $\left(x-7\right)^2+\left(y-6\right)^2+\left(z+5\right)^2=81$}
{$\left(x-7\right)^2+\left(y-6\right)^2+\left(z+5\right)^2=9$}
\end{ex}
\begin{ex}%Câu 9
Bảng số liệu ghép nhóm về chiều cao đo được (đơn vị: cm) của $30$ học sinh nam lớp 12A2 đầu năm học $2024-2025$ của một trường THPT được cho như sau:\\
\centerline{\begin{tblr}{|c|c|c|c|c|c|}
 \hline
 Chiều cao & $\left[150;\ 155\right)$ & $\left[155;\ 160\right)$ & $\left[160;\ 165\right)$ & $\left[165;\ 170\right)$ & $\left[170;\ 175\right)$\\
 \hline
 Tần số & $ 3$ & $ 7$ & $ 10$ & $ 7$ & $ 3$\\
 \hline
\end{tblr}}\\
Tính độ lệch chuẩn của mẫu số liệu ghép nhóm trên.
\choice
{\True $\dfrac{\sqrt{285}}{3}$}
{$\dfrac{\sqrt{287}}{3}$}
{$ 4\sqrt{2}$}
{$\sqrt{71}$}
\end{ex}
\begin{ex}%Câu 10
Cho biết $\lim\limits_{x\to\sqrt{3}}\dfrac{2x^2-6}{x-\sqrt{3}}=a\sqrt{b}$ ($a$, $b$ nguyên và $b<10$). Khi đó giá trị của $P=a+b$ bằng
\choice
{\True $7$}
{$10$}
{$5$}
{$6$}
\end{ex}
\begin{ex}%Câu 11
Cho $\int\dfrac{1}{x{\ln^2}x}\mathrm{\,d}x=F(x)+C$. Khẳng định nào dưới đây đúng?
\choice
{$F'(x)=-\dfrac{1}{\ln x}$}
{$F'(x)=-\dfrac{1}{\ln x}+C$}
{\True $F'(x)=\dfrac{1}{x{\ln^2}x}$}
{$F'(x)=-\dfrac{1}{\ln^2x}$}
\end{ex}
\begin{ex}%Câu 12
\immini
{
 Diện tích phần gạch sọc trong hình vẽ bằng
\choice
{$\int\limits_{-3}^1\left|-x^2-2x-3\right|\mathrm{\,d}x$}
{$\int\limits_{-3}^1\left(x^2-2x-3\right)\mathrm{\,d}x$}
{$\int\limits_{-3}^1\left(x^2+2x-3\right)\mathrm{\,d}x$}
{\True $\int\limits_{-3}^1\left(-x^2-2x+3\right)\mathrm{\,d}x$}
}
{
 \begin{tikzpicture}[>=stealth, line join=round, line cap=round, font=\footnotesize, scale=1,thick,x=.5cm,y=.5cm,
 declare function={f(\x)=-1*(\x)+1;g(\x)=1*(\x)^2+1*(\x)-2;}]
 \draw[->] (-3.5,0)--(2.5,0) node[below left] {$x$};
 \draw[->] (0,-3)--(0,4.5) node[below left] {$y$};
 \draw (0,0) node [below left] {$O$};
 \foreach \x/\nx in {-3/-3,1/1}
 \draw (\x,1pt)--(\x,-1pt) node [below] {$\nx$};
 \draw[pattern = north east lines, pattern color=green, line width = 1.2pt,draw=none] (-3,{f(3)}) plot[domain=1:-3] (\x, {g(\x)});
 \draw[dashed] (-3,0)--(-3,{f(-3)});
 \draw[samples=200,domain=-3.5:1.5,smooth] plot (\x,{f(\x)});
 \draw[samples=200,domain=-3.1:1.2,smooth,variable=\x] plot (\x,{g(\x)});
 \end{tikzpicture}
}
\end{ex}
\Closesolutionfile{ans}
\cauds
\Opensolutionfile{ans}[ans/ans-HXN-\sode-TF]
\begin{ex}%Câu 13
    Cho hàm số $f(x)=\cos 2x+2x+1$
    \choiceTF
    {\True $ f\left(\dfrac{\pi}{2}\right)=\pi $}
    {Đạo hàm của hàm số là $f'(x)=2\sin 2x+2$}
    {\True Nghiệm của phương trình $f'(x)=0$ trên đoạn $\left[-\dfrac{\pi}{2};\pi\right]$ là $ x=\dfrac{\pi}{4}$}
    {Tổng giá trị lớn nhất và giá trị nhỏ nhất của hàm số trên đoạn $\left[-\dfrac{\pi}{2};\pi\right]$ bằng $ 2\pi $}
    \loigiai{
        \begin{itemchoice}
            \itemch Ta có: $f\left(\dfrac{\pi }{2}\right)=\cos \left(2\cdot \dfrac{\pi }{2}\right)+2\cdot \dfrac{\pi }{2}+1=\pi $.
            \itemch Ta có $f'(x)=-2\sin 2x+2$.
            \itemch $f'(x)=0\Leftrightarrow -2\sin 2x+2=0\Leftrightarrow \sin 2x=1\Leftrightarrow 2x=\dfrac{\pi }{2}+k2\pi \,,k\in \mathbb{Z}\Leftrightarrow x=\dfrac{\pi }{4}+k\pi ,k\in \mathbb{Z}$.
            Vì $x\in \left[-\dfrac{\pi }{2};\pi \right]$ nên $x=\dfrac{\pi }{4}$ (ứng với $k=0$).
            \itemch Ta có: $f\left(-\dfrac{\pi }{2}\right)=-\pi $; $f\left(\dfrac{\pi }{4}\right)=\dfrac{\pi }{2}+1$; $f\left(\pi \right)=2\pi +2$.\\
            Vì vậy $\max\limits_{\left[-\dfrac{\pi}{2};\pi \right]} \,f(x)=f\left(\pi \right)=2\pi +2$ và $\min\limits_{\left[-\dfrac{\pi}{2};\pi \right]}\,f(x)=f\left(-\dfrac{\pi }{2}\right)=-\pi $.\\
            Tổng giá trị của chúng là $2\pi +2+\left(-\pi \right)=\pi +2$.
        \end{itemchoice}
    }
\end{ex}
\begin{ex}%Câu 14
\immini
{
 Trong một hệ trục $Oxyz$ cho trước, đơn vị trên mỗi trục là mét, có hai cậu bé đang chơi bắn bi trên mặt đất phẳng (cũng là mặt phẳng $Oxy$). Gần vị trí các cậu bé có một bức tường được mô hình hóa bởi mặt phẳng có phương trình $x+y-10=0$.\\
Cậu bé An bắn viên bi từ vị trí $ A\left(0;0;0\right)$ và theo hướng vectơ $\vec{u}=\left(4;-1;0\right)$, bi có thể di chuyển $10{,}5$ mét nếu không có vật cản.
}
{
 \includegraphics[width=5.5cm]{img/HXN-12-14}
}
Cậu bé Cò bắn viên bi từ vị trí $B\left(2;-2;0\right)$ theo hướng $\vec{v}=\left(1;2;0\right)$, bi có thể di chuyển $8{,}2$ mét nếu không có vật cản.
\choiceTF
 {\True Đường đi viên bi là của An có phương trình $\heva{& x=4t\\& y=-t\\& z=0}$ ($t$ là tham số)}
{\True Đường đi của hai viên bi là các đường thẳng cắt nhau tại điểm có tung độ âm}
 {Đường đi của viên bi do bạn Cò bắn ra tạo với bức tường một góc thuộc khoảng $\left(69^{\circ};71^{\circ}\right)$}
 {\True Một trong hai viên bi chạm vào bức tường}
 \loigiai{
 \begin{itemchoice}
 \itemch Đường đi viên bi là của An có phương trình $d_1\colon \heva{& x=4t \\& y=-t \\& z=0 } $ ($t$ là tham số).
 \itemch Đường đi viên bi là của Cò có phương trình $d_2\colon \heva{& x=2+t' \\& y=-2+2t' \\& z=0 } $ ($t$ là tham số).\\
 Xét hệ phương trình $\heva{& 4t=2+t' \\& -t=-2+2t' \\& 0=0 } \Rightarrow \heva{& t=\dfrac{2}{3} \\& t'=\dfrac{2}{3} } $; đường đi $2$ viên bi cắt nhau tại $M\left(\dfrac{8}{3};-\dfrac{2}{3};0\right)$; trong đó $y_M=-\dfrac{2}{3}<0$ (thỏa mãn).
 \itemch Đường thẳng $d_2$ có vectơ chỉ phương $\vec{v}=(1;2;0)$; mặt phẳng chứa bức tường có vectơ pháp tuyến $\vec{n}=(1;1;0)$ nên $\sin \left(d_2\,,(P)\right)=\dfrac{\left| \vec{v}\cdot \vec{n} \right|}{\left| {\vec{v}} \right|\cdot \left| {\vec{n}} \right|}=\dfrac{|1+2|}{\sqrt{5}\cdot \sqrt{2}}=\dfrac{3\sqrt{10}}{10}$.\\
 $\Rightarrow \left(d_2,(P)\right)\approx 71,6^{\circ }\notin \left(69^{\circ };71^{\circ }\right)$.\\
 \itemch Xét $d_1$ và $(P)\colon x+y-10=0$;\\
 thay phương trình $d_1$ vào $(P)$, ta có $4t+(-t)-10=0\Rightarrow t=\dfrac{10}{3}$.\\
 Do đó $d_1\cap (P)=C\left(\dfrac{40}{3};-\dfrac{10}{3};0\right)$ với $AC=\dfrac{10\sqrt{17}}{3}\approx 13{,}7\,m>10{,}5\,m$.\\
 Vậy viên bi của An không chạm tường.\\
 Xét $d_2$ và $(P)\colon x+y-10=0$;\\
 thay phương trình $d_2$ vào $(P)$, ta có $(2+t)+\left(-2+2t'\right)-10=0\Rightarrow t'=\dfrac{10}{3}$.\\
 Do đó $d_2\cap (P)=D\left(\dfrac{16}{3};\dfrac{14}{3};0\right)$ với $BD=\dfrac{10\sqrt{5}}{3}\approx 7{,}45\,m<8{,}2\,m$.\\
 Vậy viên bi của Cò chạm vào tường.
 \end{itemchoice}
 }
\end{ex}
\begin{ex}%Câu 15
\immini
{
 Trong một live show âm nhạc có ca sĩ Mỹ Tâm tham gia, nhiều fan hâm mộ đã tỏ ra lo ngại rằng ban tổ chức có thể hủy show với một vài lý do khác nhau. Những lo ngại này là đúng vì có đến $0{,}302$ khả năng show diễn sẽ bị hủy.\\
Nếu vé bán hết thì chắc chắn live show sẽ diễn ra.\\
Nếu trời mưa thì ban tổ chức không thể bán hết vé, khi đó khả năng hủy show là $50\%$.\\
Nếu trời không mưa thì khả năng vé được bán hết là $90\%$; trong trường hợp còn vé thì khả năng hủy show là $5\%$.
}
{
 \includegraphics[width=5.5cm]{img/HXN-12-15}
}
\choiceTF
 {\True Nếu vé chưa được bán hết khi trời không mưa thì xác suất để show vẫn diễn ra bằng $0{,}95$}
 {Xác suất để trời mưa bằng $0{,}55$}
 {\True Xác suất để ban tổ chức không bán hết vé bằng $0{,}64$}
 {\True Sau cùng thì show của Mỹ Tâm cũng đã diễn ra, xác suất hôm đó trời mưa bằng $0{,}43$ (làm tròn kết quả đến hàng phần trăm)}
 \loigiai{
 \begin{itemchoice}
 \immini
 {
 \itemch Nếu vé chưa được bán hết khi trời không mưa thì xác suất để show vẫn diễn ra bằng $1-0{,}05=0{,}95$.
 \itemch Gọi A là biến cố \lq\lq Trời mưa\rq\rq, B là biến cố: \lq\lq Vé được bán hết\rq\rq, C là biến cố: \lq\lq Show bị hủy\rq\rq.\\
 Đặt $P(A)=x\in (0;1)$ là xác suất để trời mưa, ta có sơ đồ hình cây bên cạnh.\\
 Khi đó $P(C)=x\cdot 1.0{,}5+(1-x)\cdot 0{,}9.0+(1-x)\cdot 0{,}1.0{,}05=0{,}302\Rightarrow x=0{,}6$ hay $P(A)=0{,}6$.
 }
 {
 \includegraphics[width=5.5cm]{img/HXN-12-15-LG}
 }
 \itemch Ta có $P\left({\bar{B}}\right)=\underbrace{x\cdot 1+(1-x)\cdot 0{,}1}_{x=0{,}6}=0{,}64$.
 \itemch Ta có $P\left({\bar{C}}\right)=1-0{,}302=0{,}698$; $P\left(A|\bar{C}\right)=\dfrac{P\left(A\bar{C}\right)}{P\left({\bar{C}}\right)}=\dfrac{0{,}6.1\cdot 0{,}5}{0{,}698}=\dfrac{150}{349}\approx 0{,}43$.
 \end{itemchoice}
 }
\end{ex}
\begin{ex}%Câu 16
    \immini
    {
        Cho đường cong $(C)\colon y=8x-27x^3$ và đường thẳng $d\colon y=m$ cắt $(C)$ tại ba điểm phân biệt, trong đó có hai điểm nằm ở góc phần tư thứ nhất của hệ tọa độ $ Oxy$. Các phần diện tích tô đậm được kí hiệu là $S_1$, $S_2$ như hình vẽ.
        \choiceTF
        {Diện tích hình phẳng giới hạn vởi $(C)$ và $Ox$ bằng $\dfrac{34}{27}$}
        {Có ba giá trị nguyên của $m$ thỏa mãn đề bài}
        {\True Nếu $ m=\dfrac{5}{8}$ thì $S_1+S_2=0,31$ (làm tròn đến hàng phần trăm)}
        {Nếu $S_1=S_2$ thì có một giá trị $ m=m_0$ thỏa mãn với $ 1<m_0<\dfrac{3}{2}$}
    }
    {
        \begin{tikzpicture}[>=stealth, line join=round, line cap=round, thick,declare function={f(\x)=-27*((\x)^3)+0*((\x)^2)+8*(\x)+0;m=7/8;g(\x)=m;},xscale=3]
            \draw[->] (-1,0)--(1,0) node[below left] {$x$};
            \draw[->] (0,-2.5)--(0,2.5) node[below left] {$y$};
            \draw (0,0) node [below left] {$O$};
            \begin{scope}
                \clip (-1,-2.2) rectangle (1,2.2);
                \draw[samples=200,domain=-2:2,smooth,name path=df] plot (\x,{f(\x)});
                \draw[name path=dg] (-2,m)--(2,m)node[pos=.72,above]{$d$};
                \path [name intersections= {of = df and dg, by={a,b,c}}];
            \end{scope}
            \fill[pattern=north east lines] let \p2=(b),\n2={\x2/1cm} in plot[domain=0:\n2](\x,{f(\x)})--(0,m)--cycle;
            \fill[pattern=north west lines,pattern color=red] let \p2=(b),\n2={\x2/1cm},\p3=(c),\n3={\x3/1cm} in plot[domain=\n2:\n3](\x,{f(\x)})--cycle;
            \draw[->] (.07,.7)--++(-155:.3)node[below left]{$S_1$};
            \draw[->] (.3,1.3)--++(55:.5)node[above right]{$S_2$};
        \end{tikzpicture}
    }
    \loigiai{
        \begin{itemchoice}
            \itemch Hoành độ giao điểm của $(C)$ và $Ox$ thỏa mãn $8x-27x^3=0\Leftrightarrow \hoac{& x=0 \\& x=\pm \dfrac{2\sqrt{6}}{9} } $.\\
            Diện tích cần tính là $S=\displaystyle\int\limits_{-\tfrac{2\sqrt{6}}{9}}^{\tfrac{2\sqrt{6}}{9}}{\left| 8x-27x^3 \right|\mathrm{\,d}x}=\dfrac{32}{27}$.
            \itemch Xét $(C):y=8x-27x^3$ có đạo hàm $y'=8-81x^2;y'=0\Rightarrow x=\pm \dfrac{2\sqrt{2}}{9}$.\\
            Bảng biến thiên:
            \begin{center}
                \begin{tikzpicture}[>=stealth]
                    \tkzTabInit[nocadre=false,lgt=1,espcl=2.5,deltacl=0.5]{$x$/1,$y'$/.7,$y$/2}
                    {$-\infty$ , $-\dfrac{2\sqrt{2}}{9}$ , $\dfrac{2\sqrt{2}}{9}$ , $+\infty$}
                    \tkzTabLine{ ,-,$0$,+,$0$,-, }
                    \tkzTabVar{+/$+\infty$,-/$\approx -1{,}7$,+/$\approx 1{,}7$,-/$-\infty$}
                \end{tikzpicture}
            \end{center}
            Xét đường thẳng d:$y=m$ cắt đồ thị hàm số tại ba điểm phân biệt (giả thiết); mà $m$ nguyên nên $m\in \left\{ -1;0;1 \right\}$. Tuy nhiên chỉ có $m=1$ thỏa mãn vì khi đó d và (C) có hai giao điểm thuộc góc phần tư thứ nhất. 
            \itemch 
            Nếu $m=\dfrac{5}{8}$ thì ta có phương trình hoành độ giao điểm: $8x-27x^3=\dfrac{5}{8}\Leftrightarrow \hoac{& x=\dfrac{1}{2} \\& x=\dfrac{-9\pm \sqrt{141}}{36} } $.\\
            Ta chỉ xét góc phần tư thứ nhất nên nhận $x=\dfrac{1}{2};x=\dfrac{-9+\sqrt{141}}{36}$.\\
            Khi đó $S_1+S_2=\displaystyle\int\limits_0^{\tfrac{1}{2}}{\left| \left(8x-27x^3\right)-\dfrac{5}{8} \right|\mathrm{\,d}x}\approx 0{,}31$.
            \itemch 
            Giả sử đường thẳng $y=m$ cắt đồ thị $(C)$ tại điểm có hoành độ $a,b\,\left(0<a<b\right)$.\\
            Ta có: $8a-27a^3=m$ \tagEX{1}
            Xét hàm số $f(x)=8x-27x^3-m$ có $F(x)=\displaystyle\int{f(x)}\mathrm{\,d}x=4x^2-\dfrac{27}{4}x^4-mx+C$.\\
            Ta có 
            $
            \begin{aligned}[t]
                &S_1=\displaystyle\int\limits_0^a{\left| f(x) \right|}\mathrm{\,d}x=-\displaystyle\int\limits_0^a{f(x)}\mathrm{\,d}x=F(0)-F(a);\\ &S_2=\displaystyle\int\limits_a^b{\left| f(x) \right|}\mathrm{\,d}x=\displaystyle\int\limits_a^b{f(x)}\mathrm{\,d}x=F(b)-F(a).
            \end{aligned}
            $\\
            Theo giả thiết: $S_1=S_2\Leftrightarrow F(0)-F(a)=F(b)-F(a)\Leftrightarrow F(b)=F(0)$.\\
            Vì vậy $4b^2-\dfrac{27}{4}b^4-mb=0$ \tagEX{2}
            Từ $(1)$ và $(2)$ ta có $\heva{& 4b^2-\dfrac{27b^4}{4}-mb=0 \\& 8b-27b^3=m } \Rightarrow \heva{& 4b^2-\dfrac{27b^4}{4}-\left(8b-27b^3\right)b=0 \\& 8b-27b^3=m } \Rightarrow \heva{& b=\dfrac{4}{9}>0 \\& m=\dfrac{32}{27}.} $\\
            Vậy $m=\dfrac{32}{27}\approx 1{,}19\in \left(1;\dfrac{3}{2}\right)$.
        \end{itemchoice}
    }
\end{ex}
\Closesolutionfile{ans}
\caukq
\Opensolutionfile{ans}[ans/ans-HXN-\sode-SA]
\begin{ex}%Câu 17
 Cho hình chóp $S.ABCD$ có đáy là hình vuông cạnh $ 1$, tam giác $SAB$ đều và nằm trong mặt phẳng vuông góc với đáy. Tìm khoảng cách giữa hai đường thẳng SA, BC và làm tròn đến hàng phần trăm.
 \shortans{0,87}
 \loigiai{
\immini
{
         Gọi $H$ là trung điểm $AB$, suy ra $SH\perp AB$ (do tam giác $SAB$ đều).\\
     Mặt khác $(SAB)\perp (ABCD)$ nên $SH\perp (ABCD)$.\\
     Ta có $\heva{& (SAB)\perp (ABCD) \\& AB=(SAB)\cap (ABCD) \\& BC\perp AB } \Rightarrow BC\perp (SAB)$ \tagEX{1}
     Trong mặt phẳng $(SAB)$, dựng $BK\perp SA$ tại $K$ \tagEX{2}
     Từ $(1)$, $(2)$ suy ra $BK$ là đoạn vuông góc chung của $SA$ và $BC$.
     Vậy $d\left(SA, BC\right)=BK=\dfrac{\sqrt{3}}{2}\approx 0{,}87$.
}
{
\includegraphics[width=5cm]{img/HXN-12-17-LG}
}
 }
 \end{ex}
 \begin{ex}%Câu 18
\immini
{
    Một trò chơi điện tử quy định như sau: Có 5 trụ $ A$, $B$, $C$, $D$, $E$ với số lượng các thử thách trên đường đi giữa các cặp trụ được mô tả trong hình bên. Người chơi xuất phát từ một trụ nào đó, đi qua tất cả các trụ còn lại, mỗi khi đi qua trụ nào thì trụ đó sẽ bị phá hủy và không thể quay trở lại trụ đó được nữa, nhưng người chơi vẫn phải trở về trụ ban đầu. Tổng số thử thách của đường đi thoả mãn điều kiện trên nhận giá trị nhỏ nhất là bao nhiêu?
\shortans{45}
}
{
    \includegraphics[width=5cm]{img/HXN-12-18}
}
\end{ex}
\begin{ex}%Câu 19
\immini
{
    Giả sử số lượng của một quần thể nấm men tại môi trường nuôi cấy trong phòng thí nghiệm được mô hình hóa bằng hàm số $ P(t)=\dfrac{a}{b+e^{-0,75t}}$. Trong đó, thời gian $ t$ được tính bằng giờ. Tại thời điểm ban đầu $t=0$, quần thể có 20 tế bào và tăng với tốc độ $12$ tế bào/giờ. Số lượng của quần thể nấm này tại thời điểm $ t=8$ giờ là bao nhiêu tế bào (làm tròn đến hàng đơn vị)?
\shortans{99}
}
{
    \includegraphics[width=5cm]{img/HXN-12-19}
}
\loigiai{
    Ta có $P'(t)=\dfrac{0{,}75a{e^{-0{,}75t}}}{\left(b+{e^{-0{,}75t}}\right)^2}$, $t\ge 0$.\\
    Theo giả thiết: $\heva{& P(0)=20 \\& P'(0)=12 }\Leftrightarrow \heva{& \dfrac{a}{b+1}=20 \\& \dfrac{0{,}75a}{(b+1)^2}=12 } \Leftrightarrow \heva{& a=20(b+1) \\& \dfrac{0{,}75\times 20(b+1)}{(b+1)^2}=12 } \Leftrightarrow \heva{& a=25 \\& b=\dfrac{1}{4}.}$\\
    Do vậy $P(t)=\dfrac{25}{\dfrac{1}{4}+{e^{-0.75t}}}\Rightarrow P(8)\approx 99$ (tế bào).
}
\end{ex}
\begin{ex}%Câu 20
\immini
{
    Trong không gian với hệ trục $Oxyz$ thích hợp, mặt đất là mặt phẳng Oxy, đơn vị trên mỗi trục là mét, một con chim bói cá bay xuống đớp con mồi rồi bay lên khỏi mặt nước với một đường cong parabol đẹp mắt trước khi đậu vào một nhành cây. Biết rằng đường parabol này đi qua các điểm $ A\left(1;2;2\right)$, $B\left(1;3;8\right)$, $C\left(1;-1;10\right)$, vị trí chim bói cá đậu sau khi bắt được mồi có tung độ bằng 4,1 mét. Tính chiều cao của chim đang đậu so với mặt đất theo đơn vị mét và làm tròn đến hàng phần chục.
\shortans{19,6}
}
{
    \includegraphics[width=5cm]{img/HXN-12-20}
}
\loigiai{
    Ta nhận thấy cả ba điểm $A(1;2;2)$, $B(1;3;8)$, $C(1;-1;10)$ cùng thuộc mặt phẳng $x=1$.\\
    Do vậy đường parabol mà con chim bói cá này vẽ nên thuộc mặt phẳng $x=1$.\\
    Đặt $z=ay^2+by+c$, $\left(a\ne 0\right)$ là hàm số bậc hai mô phỏng parabol $(P)$ thuộc mặt phẳng $x=1$.\\
    Thay các cặp $(y;z)$ gồm $(2;2)$, $(3;8)$, $C(-1;10)$ vào hàm số trên ta được hệ phương trình \\
    \centerline{
        $\heva{& 4a+2b+c=2 \\& 9a+3b+c=8 \\& a-b+c=10 } \Leftrightarrow \heva{& a=\dfrac{13}{6} \\& b=-\dfrac{29}{6} \\& c=3 } $ hay $z=\dfrac{13}{6}y^2-\dfrac{29}{6}y+3$.
    }
    Vị trí chim bói cá đậu có $y=4{,}1\Rightarrow z=\dfrac{3\,921}{200}\approx 19{,}6\,m$.\\
    Vậy vị trí chim bói cá đậu cách mặt đất khoảng $19{,}6$ m.
}
\end{ex}
\begin{ex}%Câu 21
\immini
{
    Một cái bục bằng gỗ dùng để đặt đồ trang trí có mặt đáy trên và mặt đáy dưới đều là hình vuông, người thợ thiết kế cái bục này theo ba phần:\\
Phần trên cùng là một hình hộp chữ nhật có các kích thước là $1$ m; $1$ m; $0{,}05$ m.\\
Phần đế của bục cũng là hình hộp chữ nhật có các kích thước là $\sqrt{2}m$; $\sqrt{2}m$; $0{,}2m$.
}
{
    \includegraphics[width=5cm]{img/HXN-12-21}
}
Phần thân (giữ của bục có mặt cắt theo hai đường chéo của đáy trên và đáy dưới là đường hypebol mà các đường tiệm cận của hypebol này tạo với trục đứng một góc bằng $30^{\circ}$.\\
Biết rằng mặt cắt của bục song song với hai đáy tại vị trí có kích thước hình vuông bé nhất bằng $0{,}5$ m. Tìm thể tích của cái bục đã cho theo đơn vị mét khối và làm tròn đến hàng phần trăm.
\shortans{2,33}
\loigiai{
\immini
{
Đặt hệ trục tọa độ như hình vẽ, gọi phương trình chính tắc hypebol $(H)$ là $\dfrac{x^2}{a^2}-\dfrac{y^2}{b^2}=1$, $\left(a>0,\,b>0\right)$.\\
Đường chéo hình vuông nhỏ nhất là $2a=0{,}5\sqrt{2}\Rightarrow a=\dfrac{\sqrt{2}}{4}$.\\
Ta có $\tan 30^{\circ }=\dfrac{a}{b}=\dfrac{\dfrac{\sqrt{2}}{4}}{b}\Rightarrow b=\dfrac{\sqrt{6}}{4}$.\\
Do vậy $(H)\colon \dfrac{x^2}{1/8}-\dfrac{y^2}{3/8}=1$ hay $8x^2=1+\dfrac{8y^2}{3}\Rightarrow x=\pm \sqrt{\dfrac{1}{8}+\dfrac{y^2}{3}}$.\\
Xét mặt cắt vuông góc với $Oy$ của cái bục tại vị trí có hoành độ $x>0$ thì ta thu được hình vuông có cạnh $\dfrac{2x}{\sqrt{2}}=x\sqrt{2}=\sqrt{\dfrac{1}{4}+\dfrac{2y^2}{3}}$.
}
{
    \includegraphics[width=4.6cm]{img/HXN-12-21-LG}
}
Với $x=\dfrac{1\cdot \sqrt{2}}{2}$ (xét hình vuông mặt trên cùng của cái bục); ta có $y=\dfrac{3\sqrt{2}}{4}>0$.\\
Với $x=\dfrac{\sqrt{2}\cdot \sqrt{2}}{2}=1$ (xét hình vuông mặt đáy dưới của cái bục); ta có $y=-\dfrac{\sqrt{42}}{4}<0$.\\
Thể tích cái bục là $V=1\cdot 1\cdot 0{,}05+\sqrt{2}\cdot \sqrt{2}\cdot 0{,}2+\int\limits_{\tfrac{-\sqrt{42}}{4}}^{\tfrac{3\sqrt{2}}{4}}{\left(\dfrac{1}{4}+\dfrac{2y^2}{3}\right)dy}\approx 2{,}33\,m^3$.
}
\end{ex}
\begin{ex}%Câu 22
\immini
{
Lớp 12A có tất cả $30$ học sinh, trong đó bạn An muốn làm lớp trưởng, bạn Bảo muốn làm lớp phó học tập và bạn Châu muốn làm bí thư (các học sinh của lớp không trùng tên). Tính xác suất để cô giáo chọn được học sinh của lớp cho ba chức vụ như trên, mỗi học sinh giữ một chức vụ, đồng thời không có ai trong các bạn An, Bảo, Châu được giữ chức vụ mình thích (làm tròn đến hàng phần chục)?
\shortans{0,9}
}
{
    \includegraphics[width=5cm]{img/HXN-12-22}
}
\loigiai{
    Gọi $A$ là tập hợp các cách chọn chức vụ sao cho An được làm lớp trưởng.\\
    Gọi $B$ là tập hợp các cách chọn chức vụ sao cho Bảo được làm lớp phó học tập.\\
    Gọi $C$ là tập hợp các cách chọn chức vụ sao cho Châu được làm bí thư.\\
    Số cách chọn chức vụ tùy ý: $n\left(\Omega \right)=30\cdot 29\cdot 28=24\,360$.\\
    Ta có $\mid A\cup B\cup C\mid =\mid A\mid +\mid B\mid +\mid C\mid -\mid A\cap B\mid -\mid B\cap C\mid -\mid C\cap A\mid +\mid A\cap B\cap C\mid $.\\
    Trong đó $\mid A\mid =\mid B\mid =\mid C\mid =1\cdot 29\cdot 28$ (chọn một học sinh tương ứng cho chức vụ mình thích, hai chức còn lại thì chọn tùy ý).\\
    Tương tự $\mid A\cap B\mid =\mid B\cap C\mid =\mid C\cap A\mid =1\cdot 1\cdot 28$ (chọn hai học sinh giữ chức vụ minh thích, chức còn lại thì chọn tùy ý).\\ 
    Ta có $\mid A\cap B\cap C\mid =1\cdot 1\cdot 1$ (cả ba học sinh giữ chức vụ mình thích).\\
    Vậy tổng số cách chọn thỏa mãn đề bài là: $n(X)=30\cdot 29\cdot 28-\left(3\cdot 29\cdot 28-3\cdot 28+1\right)=22\,007$.\\
    Xác suất cần tính là $P(X)=\dfrac{n(X)}{n\left(\Omega \right)}=\dfrac{22\,007}{24\,360}\approx 0{,}9$.
}
\end{ex}
\Closesolutionfile{ans}
\inputansbox{6,4,3}{ans/ans-HXN-\sode-T,ans/ans-HXN-\sode-TF,ans/ans-HXN-\sode-SA}